% Mathematical Compendium: DSP Techniques in NeonSynth
% A supplement to "The DSP Behind NeonSynth".
%
% Build: docs/theory/build_compendium.sh [out-dir]   (pdflatex, run twice for the page references)
% The version below must match project(NeonSynth VERSION ...) in CMakeLists.txt; the build script checks it.
\documentclass[11pt,letterpaper]{article}

\newcommand{\NSVersion}{1.9.0}
\newcommand{\NSDate}{17 November 2026}

% ---------------------------------------------------------------- fonts and page
\usepackage[T1]{fontenc}
\usepackage{mathptmx}                 % Times text and maths
\usepackage[scaled=0.92]{helvet}      % Helvetica for headings, header and footer
\usepackage{courier}                  % Courier for file and identifier names
\usepackage{amsmath,amssymb}
\usepackage[letterpaper,left=0.8in,right=0.8in,top=1.0in,bottom=0.95in,
            headheight=14pt,headsep=26pt,footskip=34pt]{geometry}
\usepackage{xcolor}
\usepackage{fancyhdr}
\usepackage{titlesec}
\usepackage{booktabs}
\usepackage{array}
\usepackage{enumitem}
\usepackage{tcolorbox}
\usepackage{lastpage}
\usepackage[hidelinks]{hyperref}

\hypersetup{
  pdftitle={Mathematical Compendium: DSP Techniques in NeonSynth v\NSVersion},
  pdfsubject={A Supplement to The DSP Behind NeonSynth},
}

% ---------------------------------------------------------------- colours
\definecolor{navy}{RGB}{28,42,74}         % title band, headings
\definecolor{accent}{RGB}{42,108,214}     % numbers, links, rules, box bars
\definecolor{kicker}{RGB}{150,184,238}    % title band kicker line
\definecolor{subtle}{RGB}{110,118,130}    % header text, contents descriptions
\definecolor{codegray}{RGB}{85,85,85}     % file and identifier names
\definecolor{boxbg}{RGB}{242,246,252}     % callout background
\definecolor{rulegray}{RGB}{200,204,210}

\newcommand{\code}[1]{\texttt{\textcolor{codegray}{#1}}}

% ---------------------------------------------------------------- running header and footer
\pagestyle{fancy}
\fancyhf{}
\fancyhead[L]{\sffamily\footnotesize\textcolor{navy}{\textbf{NEONSYNTH v\NSVersion}}\textcolor{subtle}{~---~MATHEMATICAL COMPENDIUM}}
\fancyhead[R]{\sffamily\footnotesize\textcolor{subtle}{DSP REFERENCE SUPPLEMENT}}
\fancyfoot[C]{\sffamily\footnotesize\textcolor{subtle}{Page \thepage{} of \textcolor{accent}{\pageref*{LastPage}}}}
\renewcommand{\headrule}{\hbox to\headwidth{\color{navy!70}\leaders\hrule height 0.6pt\hfill}}
\renewcommand{\footrule}{\hbox to\headwidth{\color{rulegray}\leaders\hrule height 0.4pt\hfill}}
\renewcommand{\footrulewidth}{0.4pt}

\fancypagestyle{titlepage}{%
  \fancyhf{}%
  \renewcommand{\headrule}{}%
  \renewcommand{\footrule}{}%
  \fancyfoot[C]{\sffamily\footnotesize\textcolor{subtle}{Target Instrument: NeonSynth v\NSVersion
    \quad|\quad Companion Reference to \textit{The DSP Behind NeonSynth}}}%
}

% ---------------------------------------------------------------- headings
\titleformat{\section}{\sffamily\bfseries\Large\color{navy}}{\textcolor{accent}{\thesection.}}{0.45em}{}
  [\vspace{0.05ex}{\color{accent}\titlerule[0.9pt]}]
\titlespacing*{\section}{0pt}{0pt}{1.8ex}
\titleformat{\subsection}{\sffamily\bfseries\large\color{navy}}{\textcolor{accent}{\thesubsection}}{0.5em}{}
\titlespacing*{\subsection}{0pt}{2.6ex}{1.3ex}

\setlength{\parindent}{0pt}
\setlength{\emergencystretch}{1.5em}
\setlength{\parskip}{0.55ex}
\setlength{\abovedisplayskip}{6pt plus 2pt minus 2pt}
\setlength{\belowdisplayskip}{6pt plus 2pt minus 2pt}
\setlength{\abovedisplayshortskip}{4pt plus 2pt}
\setlength{\belowdisplayshortskip}{6pt plus 2pt minus 2pt}
\renewcommand{\arraystretch}{1.15}

% ---------------------------------------------------------------- callouts
\newtcolorbox{callout}[1]{%
  colback=boxbg, colframe=accent, coltitle=accent, colbacktitle=boxbg,
  boxrule=0pt, leftrule=3.5pt, arc=0pt, outer arc=0pt,
  left=10pt, right=8pt, top=2pt, bottom=6pt, toptitle=5pt, bottomtitle=2pt,
  titlerule=0pt, fonttitle=\sffamily\bfseries\small,
  title={#1}, before skip=6pt, after skip=6pt}

% ---------------------------------------------------------------- notation
\newcommand{\Fs}{F_s}
\newcommand{\frc}{\operatorname{frac}}
\newcommand{\sat}{\operatorname{sat}}
\newcommand{\clamp}{\operatorname{clamp}}
\newcommand{\lfo}{\ell}

% ---------------------------------------------------------------- contents rows
\newcommand{\contentsrow}[4]{%
  \textsf{\bfseries\textcolor{accent}{#1.}} &
  \hyperref[#4]{\textbf{\textcolor{accent}{#2}}} & \\
  & {\small\textcolor{subtle}{#3}} &
  \hyperref[#4]{\sffamily\small\textcolor{subtle}{Page} \textcolor{accent}{\pageref*{#4}}}\\[0pt]
  \midrule[0.5pt]}

\begin{document}

% ================================================================ title page
\thispagestyle{titlepage}
\vspace*{0.1in}

\begin{tcolorbox}[colback=navy, colframe=accent, boxrule=0pt, bottomrule=3pt,
                  arc=4pt, outer arc=4pt, left=20pt, right=20pt, top=14pt, bottom=16pt]
  {\sffamily\bfseries\footnotesize\textcolor{kicker}{TECHNICAL REFERENCE SUPPLEMENT \textbullet{} NEONSYNTH v\NSVersion}}\par
  \vspace{2pt}
  {\sffamily\bfseries\fontsize{24}{27}\selectfont\color{white}Mathematical Compendium:\\DSP Techniques in NeonSynth\par}
  \vspace{8pt}
  {\itshape\large\textcolor{white!85!navy}{A Supplement to ``The DSP Behind NeonSynth''}\par}
\end{tcolorbox}

\vspace{0.22in}

\begin{tcolorbox}[colback=boxbg, colframe=accent, coltitle=accent, colbacktitle=boxbg,
                  boxrule=0pt, leftrule=3.5pt, arc=0pt, outer arc=0pt,
                  left=10pt, right=8pt, top=3pt, bottom=7pt, toptitle=5pt, bottomtitle=2pt,
                  titlerule=0pt, fonttitle=\sffamily\bfseries\small, title={Scope \& Edition Notes}]
This compendium presents the core mathematical formulations underlying seven digital signal processing
subsystems in \textbf{NeonSynth v\NSVersion} (released \NSDate): 4-point cubic B-spline BLEP/BLAMP
anti-aliasing, trapezoidal topology-preserving state-variable filtering, quadrature van der Pol self-oscillation
limiting, double-precision antiderivative anti-aliasing (ADAA) drive, polynomial-windowed granular synthesis,
prefiltered Catmull--Rom sample resampling, and the per-sample LFO, pitch-bend and vibrato modulation new in v1.9.0.\\
All equations, boundary conditions, and numerical thresholds are aligned with the NeonSynth v\NSVersion{} implementation
(\code{Oscillator.h}, \code{CytomicSvf.h}, \code{NeonSynthVoice.h}, \code{GrainCloud.h}, \code{SampleData.h},
\code{SampleReader.h}, \code{Lfo.h}, and \code{Modulation.h}). This edition adds Section~7 and corrects several v1.8.0 figures against the code and its tests.
\end{tcolorbox}

\vspace{0.12in}
{\sffamily\bfseries\large\color{navy}Contents\par}
\vspace{2pt}
{\color{accent}\hrule height 0.8pt}
\vspace{6pt}

\noindent
\begin{tabular*}{\textwidth}{@{}p{0.45in}@{}p{5.85in}@{\extracolsep{\fill}}r@{}}
\contentsrow{1}{Band-Limited Edge Correction (BLEP \& BLAMP)}
  {Phase accumulation, waveform latching, cubic B-spline residuals, and treble compensation}{sec:blep}
\contentsrow{2}{Topology-Preserving Transform (TPT / ZDF) Filter}
  {Bilinear pre-warping, JUCE-ordered difference equations, four filter modes, and decimation}{sec:tpt}
\contentsrow{3}{Van der Pol Self-Oscillation Limiter}
  {Resonance mapping, quadrature amplitude sensing, nonlinear damping, and degree-1 homogeneity}{sec:vdp}
\contentsrow{4}{Antiderivative Anti-Aliasing (ADAA) Drive}
  {Normalized saturation, double-precision secant evaluation, cancellation guard, and dry crossfade}{sec:adaa}
\contentsrow{5}{Granular Synthesis Mathematics}
  {Polynomial window, uncorrelated overlap gain scaling, onset jitter, and 32.32 fixed-point SIMD}{sec:gran}
\contentsrow{6}{Sample Resampling and Interpolation}
  {Playback rate, Horner Catmull--Rom cubics, 95-tap Blackman FIR decimation, and transposition bounds}{sec:smp}
\contentsrow{7}{Modulation: LFO, Pitch Bend and Vibrato}
  {LFO shapes and phase, slew limiter, pitch ratio, tremolo, per-source pitch, and bit-identity}{sec:mod}
\end{tabular*}

% ================================================================ 1
\newpage
\section{Band-Limited Edge Correction (BLEP \& BLAMP)}\label{sec:blep}

To suppress aliasing from discontinuities in naive periodic waveforms, NeonSynth applies polynomial residuals
derived from a centered, unit-area 4-point cubic B-spline kernel $B_3(x)$ supported on $[-2,2]$ (\code{Oscillator.h},
\code{Blep4}).

\subsection{Phase Accumulation and Waveform Latching}

The phase accumulator retains phase across notes and advances per sample $n$ according to:
\begin{equation}\label{eq:phase}
  \phi[n+1] = \frc\!\left(\phi[n] + \Delta\phi\right), \qquad \Delta\phi = \min\!\left(\frac{f}{\Fs},\, 0.49\right)
\end{equation}
where $f$ is the oscillator frequency, $\Fs$ is the host sample rate, and the phase increment $\Delta\phi$ is strictly capped at
0.49 cycles per sample to remain below the Nyquist limit (0.5). Since v1.9.0 the per-sample pitch ratio $r[n]$
multiplies $f/\Fs$ before the cap (Section~\ref{sec:mod}, eq.~\eqref{eq:oscstep}). Calling \code{setWaveform()} stores a pending
shape that is latched unconditionally at the phase wrap $\phi \ge 1.0$. Across a waveform switch, the downstream 3-tap
treble-compensation FIR filter (active on all non-sine waveforms; the sine bypasses it but keeps its history) preserves its state history:
\begin{equation}\label{eq:treble}
  y_{\mathrm{comp}}[n] = (1+2a)\,x[n-1] - a\bigl(x[n] + x[n-2]\bigr), \qquad a = 0.27
\end{equation}

\subsection{Residual Formulation and Boundary Conditions}

Let $d \in [0,2]$ denote the unsigned fractional sample distance from a discontinuity, and let $s \in \{-1,+1\}$ indicate the
pre-edge versus post-edge sign for step discontinuities. Because a BLEP step residual is odd-symmetric across an
amplitude jump whereas a BLAMP ramp residual is even-symmetric across a slope corner, the corrected sample is
formed by adding signed step residuals and unsigned corner residuals to the naive waveform:
\begin{equation}\label{eq:residuals}
  y[n] = y_{\mathrm{naive}}[n] + \sum_{\text{edges}} s \cdot J \cdot R_{\mathrm{BLEP}}(d) + \sum_{\text{corners}} S \cdot R_{\mathrm{BLAMP}}(d)
\end{equation}
where $J$ is the amplitude jump height and $S$ is the slope change per sample.
\begin{itemize}[leftmargin=1.6em, itemsep=2pt, topsep=4pt]
  \item \textbf{BLEP (Band-Limited Step):} The first integral of the 4-point cubic B-spline kernel:
  \begin{equation}\label{eq:blep}
    R_{\mathrm{BLEP}}(d) = \int_d^2 B_3(x)\,dx, \qquad R_{\mathrm{BLEP}}(0) = \frac{1}{2}, \quad R_{\mathrm{BLEP}}(2) = 0.
  \end{equation}
  \item \textbf{BLAMP (Band-Limited Ramp):} The second integral of the cubic B-spline kernel, used for slope discontinuities
  such as triangle wave corners:
  \begin{equation}\label{eq:blamp}
    R_{\mathrm{BLAMP}}(d) = \int_d^2\!\!\int_x^2 B_3(t)\,dt\,dx, \qquad R_{\mathrm{BLAMP}}(0) = \frac{7}{30}, \quad R_{\mathrm{BLAMP}}(1) = \frac{1}{120}.
  \end{equation}
\end{itemize}

\begin{callout}{Practical Anti-Aliasing Window Bounds}
Because each residual spans $\pm 2$ samples (4 samples total), adjacent correction windows remain completely
non-overlapping for $\Delta\phi \le 0.25$ on single-edge waveforms (sawtooth) and $\Delta\phi \le 0.125$ on two-edge waveforms
(square and triangle). For higher increments up to the 0.49 cap, windows overlap and some aliasing returns,
though output amplitude remains bounded. As a finite-support spline kernel, \code{Blep4} provides practical
anti-aliasing rather than ideal brick-wall filtering.
\end{callout}

% ================================================================ 2
\newpage
\section{Topology-Preserving Transform (TPT / ZDF) Filter}\label{sec:tpt}

NeonSynth v\NSVersion{} implements a 2-pole (12\,dB/oct) multi-mode state-variable filter per voice (\code{CytomicSvf.h})
using trapezoidal integration to eliminate the artificial unit-sample delay ($z^{-1}$) introduced by naive bilinear feedback
loops.

\subsection{Continuous-to-Discrete Mapping and Cutoff Modulation}

The continuous integrator $1/s$ is mapped to discrete time via the bilinear transform with frequency pre-warping at
cutoff $f_c$ (evaluated in 64-bit \code{double} precision and cast to \code{float}):
\begin{equation}\label{eq:tpt}
  g = \tan\!\left(\pi\,\frac{f_c}{\Fs}\right), \qquad R_2 = k_{\mathrm{eff}}, \qquad h = \frac{1}{1 + R_2\,g + g^2}
\end{equation}
At neutral settings (0\% resonance, $Q = 1/\sqrt{2}$ Butterworth, and 0\% drive), $R_2 = k = 1/Q = \sqrt{2}$ and $h$ is
evaluated in \code{double}, as JUCE does. At any resonance above 0\% the voice runs the nonlinear step, in which $R_2$ is the
amplitude-limited effective damping $k_{\mathrm{eff}}$ of Section~\ref{sec:vdp}, re-evaluated every sample (with $h$ in \code{float}).

Filter cutoff is modulated exponentially in octaves by a dedicated filter \code{VoiceEnvelope} (defaults: 10\,ms Attack,
300\,ms Decay, 0.5 Sustain, 300\,ms Release) and note velocity, clamped safely below Nyquist:
\begin{equation}\label{eq:cutoff}
  f_c = \clamp\!\left(f_{\mathrm{knob}} \cdot 2^{\,\mathrm{envAmount}\cdot\mathrm{envLevel} + \mathrm{velAmount}\cdot\mathrm{velocity}},\; 20\ \mathrm{Hz},\; 0.45\Fs\right)
\end{equation}
with envAmount $\in[-5,5]$ octaves, velAmount $\in[0,4]$ octaves and velocity $\in[0,1]$; v1.9.0 adds the LFO's cutoff octaves
to the exponent (eq.~\eqref{eq:modcutoff}). When the summed exponent is exactly zero, $f_{\mathrm{knob}}$ is returned untouched without
evaluating $2^0$. To reduce trigonometric
overhead during sweeps, coefficients $(g, R_2, h)$ update every 16 samples while parameters, envelopes or the LFO are moving
(persisting the 16-sample counter across buffer boundaries for block-size independence) and are skipped at rest.

\subsection{Sample-Step Difference Equations and Filter Modes}

For an input sample $x$ and previous integrator states $s_1, s_2$, the zero-delay outputs and trapezoidal state updates
execute in the following exact sequential order:
\begin{equation}\label{eq:svf}
\begin{aligned}
  y_{\mathrm{HP}} &= h \cdot \bigl(x - s_1(g + R_2) - s_2\bigr)\\
  y_{\mathrm{BP}} &= g \cdot y_{\mathrm{HP}} + s_1\\
  s_1 &\leftarrow g \cdot y_{\mathrm{HP}} + y_{\mathrm{BP}}\\
  y_{\mathrm{LP}} &= g \cdot y_{\mathrm{BP}} + s_2\\
  s_2 &\leftarrow g \cdot y_{\mathrm{BP}} + y_{\mathrm{LP}}
\end{aligned}
\end{equation}
All four filter modes are obtained directly from these states: Low-Pass ($y_{\mathrm{LP}}$), High-Pass ($y_{\mathrm{HP}}$), Band-Pass ($y_{\mathrm{BP}}$), and
Notch ($y_{\mathrm{LP}} + y_{\mathrm{HP}}$). Mode switches crossfade over 10\,ms to prevent clicks.

\begin{callout}{IEEE-754 Associativity and Bit-Identical Compatibility}
The five-step arithmetic sequence above copies JUCE's \code{StateVariableTPTFilter} instruction for
instruction rather than Andrew Simper's algebraically equivalent reordered Cytomic equations. Because
IEEE-754 floating-point addition and multiplication are non-associative, altering the evaluation order shifts
the least significant bits of rendered audio. Preserving JUCE's exact operation order guarantees bit-identical
output against legacy v1.7.1 sessions at neutral settings (0\% resonance, 0\% drive).
\end{callout}

% ================================================================ 3
\newpage
\section{Van der Pol Self-Oscillation Limiter}\label{sec:vdp}

The Resonance control maps percentages (0\% to 100\%) to the linear damping factor $k = 1/Q$. At 98\% resonance
($k = 0$), the linear filter poles lie on the unit circle; when resonance exceeds 98\% ($k < 0$, reaching $k = -0.0275$
at 100\%), the poles move outside the unit circle and cause exponential growth unless bounded by a restoring
nonlinearity. Rather than hard-clipping an internal state, \code{CytomicSvf.h} applies smooth amplitude-dependent
van der Pol damping.

\begin{table}[h]
\centering
\small
\begin{tabular*}{\textwidth}{@{}l@{\extracolsep{\fill}}cccccccc@{}}
\toprule
\textbf{Resonance} & 0\% & 25\% & 50\% & 75\% & 90\% & 95\% & 98\% & 100\% \\
\midrule
\textbf{Damping} $k = 1/Q$ & 1.4142 & 0.8398 & 0.5000 & 0.2568 & 0.1000 & 0.0393 & 0.0000 & $-0.0275$ \\
\textbf{Quality} $Q$       & 0.71   & 1.19   & 2.00   & 3.89   & 10.0   & 25.5   & $\infty$ & $-36.3$ \\
\bottomrule
\end{tabular*}
\caption{NeonSynth v\NSVersion{} mapping from Resonance percentage to damping factor $k$ and quality factor $Q$.}
\end{table}
\vspace{-2ex}

\subsection{Quadrature Amplitude Sensing}

Near the cutoff frequency $f_c$, the two integrator states $s_1$ and $s_2$ have equal magnitude and sit approximately $90^\circ$ (a
quarter cycle) out of phase. Their squared sum provides a phase-ripple-free measure of the instantaneous squared
envelope $A^2$, avoiding second-harmonic distortion inside the damping loop and remaining stable near $f_c = \Fs/4$:
\begin{equation}\label{eq:amp}
  A^2 = \frac{s_1^2 + s_2^2}{4}
\end{equation}

\subsection{Nonlinear Damping Equation}

Effective damping $k_{\mathrm{eff}}$ increases smoothly with the ratio of current oscillation amplitude $A$ to an envelope-tracked
dynamic ceiling $L$:
\begin{equation}\label{eq:ceiling}
  L = \max\!\left(\mathrm{velocityLevel} \times \mathrm{ampEnv},\; 10^{-6}\right), \qquad \mathrm{velocityLevel} = 0.25 \times \mathrm{velocity}
\end{equation}
\begin{equation}\label{eq:keff}
  q = c\,\frac{A^2}{L^2}, \qquad k_{\mathrm{eff}} = k + \frac{B \cdot q}{B + q}, \qquad (c = 0.1,\ B = 2)
\end{equation}
As $A$ approaches $L$, $k_{\mathrm{eff}}$ rises through zero and becomes positive, stabilizing self-oscillation into a limit cycle. At
$\Fs = 48$\,kHz, the oscillation settles at a peak of approximately $1.05L$ (staying within $\pm 1$\,dB from 100\,Hz to 5\,kHz,
$-2.1$\,dB at 10\,kHz, and $-5.1$\,dB at 15\,kHz), while tracking cutoff pitch within 0.1\%.

\subsection{Degree-1 Homogeneity}

Above the $10^{-6}$ numerical floor, $k_{\mathrm{eff}}$ depends solely on the dimensionless ratio $A/L$, making the filter response
homogeneous of degree 1:
\begin{equation}\label{eq:homog}
  f(\alpha x, \alpha L) = \alpha\, f(x, L) \qquad (\alpha > 0,\ \alpha L \ge 10^{-6})
\end{equation}
Consequently, self-oscillation tracks the amplitude envelope---swelling during Attack and decaying smoothly to
silence during Release---while preserving the physical voice signal chain (\code{ampEnv} $\rightarrow$ \code{drive} $\rightarrow$ \code{filter}).

% ================================================================ 4
\newpage
\section{Antiderivative Anti-Aliasing (ADAA) Drive}\label{sec:adaa}

To introduce warm saturation ahead of the filter loop without generating feedback intermodulation or severe aliasing,
NeonSynth v\NSVersion{} applies first-order ADAA over 0--20\,dB of Drive ($D = 10^{\mathrm{drive\_dB}/20}$).

\subsection{Normalized Saturation Function}

The underlying memoryless saturation curve is:
\begin{equation}\label{eq:sat}
  \sat(u) = \frac{u}{\sqrt{1 + u^2}}
\end{equation}
Scaling the input $x[n]$ by both the dynamic ceiling $L$ and the drive factor $D$, then normalizing by $\sat(D/4)$, preserves
degree-1 homogeneity with the filter limiter:
\begin{equation}\label{eq:drive}
  u_n = \frac{D \cdot x[n]}{4L}, \qquad y_{\mathrm{ideal}}[n] = L \cdot \frac{\sat(u_n)}{\sat(D/4)}
\end{equation}
A peak input at $|x| = L$ maintains unity peak $L$, whereas small signals ($|x| \ll L$) receive a linear gain boost of
$(D/4)/\sat(D/4) = \sqrt{1 + (D/4)^2}$, reaching $+8.6$\,dB at full 20\,dB drive ($D = 10$).

\subsection{First-Order Antiderivative Evaluation}

The continuous-time first antiderivative of $\sat(u)$ is:
\begin{equation}\label{eq:antideriv}
  F_1(u) = \int \frac{u}{\sqrt{1 + u^2}}\,du = \sqrt{1 + u^2}
\end{equation}
First-order ADAA replaces instantaneous sampling of $\sat(u)$ with its average slope across the interval $[u_{n-1}, u_n]$:
\begin{equation}\label{eq:adaa}
  y_{\mathrm{ADAA}}[n] = \frac{F_1(u_n) - F_1(u_{n-1})}{u_n - u_{n-1}}, \qquad y_{\mathrm{drive}}[n] = L \cdot \frac{y_{\mathrm{ADAA}}[n]}{\sat(D/4)}
\end{equation}

\subsection{Numerical Stability and Group-Delay Crossfade}

When successive normalized inputs are nearly equal ($\Delta u = |u_n - u_{n-1}| \le 10^{-5}$), subtracting two nearly identical
square roots causes catastrophic cancellation and risks $0/0$ division. In that regime, the implementation falls back to
the midpoint evaluation:
\begin{equation}\label{eq:midpoint}
  y_{\mathrm{ADAA}}[n] = \sat\!\left(\frac{u_n + u_{n-1}}{2}\right) \qquad \text{if } |u_n - u_{n-1}| \le 10^{-5}
\end{equation}
All $F_1(u)$ evaluations are executed in 64-bit \code{double} precision to preserve numerical accuracy near $\sqrt{1 + u^2} \approx 1$.

\begin{callout}{Half-Sample Delay Compensation and Aliasing Reduction}
Because first-order ADAA computes the average over $[u_{n-1}, u_n]$, it introduces an inherent 0.5-sample group
delay. To prevent comb filtering when engaging subtle amounts of drive, the first 5\% of the Drive knob
crossfades from the undelayed dry signal into the ADAA output: the stage outputs $x + \beta\,(y_{\mathrm{drive}} - x)$ with
$\beta = \min(1,\, 20\cdot\mathrm{Drive})$, Drive $\in[0,1]$. Measured on a PolyBLEP sawtooth at 100\%
drive, ADAA holds the aliasing floor at $-32.4$\,dB against a $-32.7$\,dB baseline, outperforming naive clipping
by 11 to 17\,dB.
\end{callout}

% ================================================================ 5
\newpage
\section{Granular Synthesis Mathematics}\label{sec:gran}

The granular engine (\code{GrainCloud.h}) overlaps short windowed excerpts of a prepared mono sample, decoupling
the source scan trajectory from per-grain playback rate:
\begin{equation}\label{eq:grain}
  \mathrm{grain}_j[n] = w\!\left(\frac{n - \mathrm{onset}_j}{\mathrm{duration}_j}\right) \times \mathrm{sample}\bigl(\mathrm{start}_j + \mathrm{rate}_j\,(n - \mathrm{onset}_j)\bigr), \qquad
  \mathrm{output}[n] = G_{\mathrm{scale}} \sum_j \mathrm{grain}_j[n]
\end{equation}
where $\mathrm{sample}(\cdot)$ is the Catmull--Rom read of eq.~\eqref{eq:horner} on the copy chosen for the grain's rate (Section~\ref{sec:smp}).

\subsection{Polynomial Window Function}

Rather than evaluating a trigonometric raised-cosine Hann window per sample, NeonSynth uses a computationally
efficient $C^1$-smooth polynomial bell on normalized grain progress $u \in [0,1]$:
\begin{equation}\label{eq:window}
  w(u) = \bigl[4u(1-u)\bigr]^2, \qquad w(0) = w(1) = 0, \quad w'(0) = w'(1) = 0, \quad w(0.5) = 1
\end{equation}
Zero value and zero slope at both endpoints eliminate boundary clicks while requiring only multiplications and
subtractions inside the inner loop.

\subsection{Overlap Gain Scaling and Pool Management}

To prevent excessive amplitude buildup as grain density or duration increases, output gain is normalized by the
square root of the expected overlap (density $\times$ duration):
\begin{equation}\label{eq:gscale}
  G_{\mathrm{scale}} = \frac{1}{\sqrt{\max(1,\ \mathrm{density} \times \mathrm{duration})}}
\end{equation}
This scaling assumes roughly uncorrelated grains. Onset intervals are jittered by $\pm 30\%$ around the nominal period to
avoid periodic ringing, and pitch jitter of $\delta$ cents scales per-grain playback rate by $2^{\delta/1200}$. Active grains are capped
at 32 simultaneous voices; when all 32 slots are occupied, new grain onsets are skipped rather than truncating active
grains.

\subsection{Fixed-Point Position Stepping and SIMD Consistency}

Grain read positions advance using 32.32 fixed-point arithmetic (32 integer bits and 32 fractional bits in a 64-bit
integer representation), enabling fast, drift-free fractional increments across long playback spans. Rendering loops
target scalar, SSE2, and AVX2 execution paths with fused multiply-add (FMA) contraction disabled so numerical
rounding remains consistent across CPU instruction sets.
Under pitch modulation (v1.9.0) the voice renders the cloud in chunks of at most 32 samples, each at the pitch ratio
of its first sample: new grains start at that pitch on the copy that suits it, and running grains rescale their integer
step once per chunk (eq.~\eqref{eq:grainstep}), before the inner loop, so the three paths stay bit-identical; at an unchanged
ratio the step is untouched.

\begin{callout}{Decoupled Scan vs.\ Per-Grain Playback}
Source selection (\code{Position} + \code{Scan} + \code{Spray}) is sampled only at each grain's onset. Freezing \code{Scan} halts
cursor movement across the source buffer while active grains continue playing to completion; negative \code{Scan}
moves the selection cursor backward through the buffer without reversing audio inside individual grains.
\end{callout}

% ================================================================ 6
\newpage
\section{Sample Resampling and Interpolation}\label{sec:smp}

Pitch-shifting sampled audio is formulated as a variable-rate discrete resampling problem (\code{SampleData.h},
\code{SampleReader.h}).

\subsection{Playback Rate Calculation}

Given a MIDI \code{note}, mapped \code{root} key, prepared file sample rate $F_{\mathrm{file}}$, and host sample rate $\Fs$, the sample read
increment per host sample is:
\begin{equation}\label{eq:rate}
  \mathrm{rate} = 2^{(\mathrm{note}-\mathrm{root})/12} \times \frac{F_{\mathrm{file}}}{\Fs}, \qquad \mathrm{position}_{n+1} = \mathrm{position}_n + \mathrm{rate}
\end{equation}
Each semitone scales the traversal rate by $2^{1/12}$, and transposing one octave upward doubles both pitch and traversal
rate. Since v1.9.0 the increment is multiplied per sample by the pitch ratio $r[n]$ (eq.~\eqref{eq:oscstep}).

\subsection{Catmull--Rom Cubic Interpolation}

At fractional sample offset $t \in [0,1)$ across four neighboring stored samples $P_0, P_1, P_2, P_3$, NeonSynth performs
Catmull--Rom cubic interpolation passing through $P_1$ and $P_2$ with $C^1$ continuity (continuous first derivative),
evaluated via Horner's nested multiplication:
\begin{equation}\label{eq:horner}
  y(t) = c_0 + t\bigl(c_1 + t(c_2 + t \cdot c_3)\bigr)
\end{equation}
where $c_0 = P_1$, $c_1 = \tfrac12(P_2 - P_0)$, $c_2 = P_0 - \tfrac52 P_1 + 2P_2 - \tfrac12 P_3$ and
$c_3 = \tfrac12(P_3 - P_0) + \tfrac32(P_1 - P_2)$ (with buffer padding ensuring boundary-safe reads).

\subsection{Prefiltered Decimation Levels}

To attenuate high-frequency content that would alias during upward transposition, \code{SampleData.h} precomputes
three versions of each imported mono sample ($-1$\,dBFS normalized, 24-bit quantized, max 60\,s) at load time:
\begin{enumerate}[leftmargin=1.6em, itemsep=0pt, topsep=2pt]
  \item Original rate ($F_{\mathrm{file}}$)
  \item Half rate ($F_{\mathrm{file}}/2$)
  \item Quarter rate ($F_{\mathrm{file}}/4$, decimated again from the half-rate table)
\end{enumerate}
The half-rate and quarter-rate tables are band-limited prior to decimation using an $N = 95$-tap Blackman-windowed
sinc low-pass FIR filter (cutoff $0.21$ of the input rate, unity gain at DC; $0 \le n \le N-1$), trading a slightly wider main lobe for strong sidelobe suppression:
\begin{equation}\label{eq:blackman}
  w[n] = 0.42 - 0.5\cos\!\left(\frac{2\pi n}{N-1}\right) + 0.08\cos\!\left(\frac{4\pi n}{N-1}\right)
\end{equation}

\begin{callout}{Transposition Bounds of the Three-Tier Table}
At note-on, \code{SampleReader.h} selects the highest-rate table whose effective playback increment satisfies
$\mathrm{rate} \le 2^{3/12} \approx 1.19$ (three semitones above unity). Because only three decimation tiers ($1\times$, $1/2\times$, $1/4\times$) are stored, extreme upward
transposition---such as playing more than 27 semitones above \code{root} when $F_{\mathrm{file}} = \Fs$ (yielding $\mathrm{rate}_{\mathrm{orig}} > 4 \cdot 2^{3/12} \approx 4.76$)---causes even the
quarter-rate table to exceed that bound. Since v1.9.0 the bound applies to $\mathrm{rate}\cdot r[n]$, and bend, LFO and vibrato can add up to 37 semitones
($r \le 2^{37/12} \approx 8.5$; eq.~\eqref{eq:levels}). This three-level architecture is therefore a practical, memory-bounded
anti-aliasing compromise rather than an arbitrary-ratio resampler.
\end{callout}

% ================================================================ 7
\newpage
\section{Modulation: LFO, Pitch Bend and Vibrato}\label{sec:mod}
\setlength{\abovedisplayskip}{4pt plus 1pt minus 1pt}\setlength{\belowdisplayskip}{4pt plus 1pt minus 1pt}
\setlength{\abovedisplayshortskip}{2pt plus 1pt}\setlength{\belowdisplayshortskip}{4pt plus 1pt minus 1pt}

One LFO (\code{Lfo.h}) is rendered once per block into $\lfo[n] \in [-1,1]$; \code{Modulation.h} adds the block's pitch bend and
mod wheel (at their MIDI sample positions) and builds per-sample buffers shared by every voice.

\subsection{LFO Shapes, Phase and Click Suppression}

On phase $p \in [0,1)$ (triangle pieces on $[0,\frac14)$, $[\frac14,\frac34)$, $[\frac34,1)$; $u_k$ uniform on $[-1,1)$ per cycle, PRNG reseeded by \code{reset()}):
\begin{equation}\label{eq:shapes}
  \lfo_{\sin} = \sin 2\pi p, \quad
  \lfo_{\mathrm{tri}} = \{4p,\ 2-4p,\ 4p-4\}, \quad
  \lfo_{\mathrm{sq}} = 1 - 2\,[p \ge \tfrac12], \quad
  \lfo_{\mathrm{saw}} = 2p - 1, \quad \lfo_{\mathrm{rnd}} = u_k
\end{equation}
\begin{equation}\label{eq:lfophase}
  p[n+1] = \frc\!\left(p[n] + \rho/\Fs\right), \quad
  \rho = \clamp(\mathrm{rate},\, 0.02,\, 20) \ \text{or}\ \frac{\mathrm{bpm}}{60B}\ \text{(synced)}, \quad
  p[n_0] = \frc(\mathrm{ppq}/B)
\end{equation}
with $B \in \{16, 8, 4, 2, 1, \frac12, \frac14, \frac13, \frac16\}$ beats (4 bars \ldots\ 1/16T); $p[n_0]$ at block start applies only when synced to a
playing host. A rate limiter follows each shape; a shape change crossfades linearly over $K = \operatorname{round}(0.01\Fs)$ samples:
\begin{equation}\label{eq:slew}
  y[n] = y[n-1] + \clamp\!\left(x[n] - y[n-1],\, \pm\Delta\right), \quad \Delta = \frac{2}{0.002\,\Fs}; \qquad
  \lfo[n] = y_{\mathrm{old}}[n] + \tfrac{k}{K}\bigl(y_{\mathrm{new}}[n] - y_{\mathrm{old}}[n]\bigr)
\end{equation}
so $|\Delta y| \le 1000/\Fs$ per sample: sine, triangle and saw ramps ($\le 126$\,s$^{-1}$ at 20\,Hz) pass bit for bit.

\subsection{Pitch Ratio, Cutoff and Tremolo}

Bend $b\in[-1,1]$ (14-bit, any channel) scaled by range $R \in 0..24$ st (default 2); wheel $w\in[0,1]$ (CC\,1, 14-bit with CC\,33; CC\,121
resets); vibrato $c \in 0..100$ cents (default 30) at 5.5\,Hz; LFO depths $d_p \in \pm 12$ st, $d_c \in \pm 5$ oct, $A\in[0,1]$:
\begin{equation}\label{eq:ratio}
  r[n] = 2^{\left(b[n]\,R \,+\, w[n]\,\frac{c}{100}\sin 2\pi\psi[n] \,+\, d_p[n]\,\lfo[n]\right)/12}, \qquad \psi[n+1] = \frc\!\left(\psi[n] + 5.5/\Fs\right)
\end{equation}
\begin{equation}\label{eq:modcutoff}
  f_c[n] = \clamp\!\left(f_{\mathrm{knob}} \cdot 2^{\,\mathrm{env} + \mathrm{vel} + d_c[n]\,\lfo[n]},\, 20\,\mathrm{Hz},\, 0.45\Fs\right), \qquad
  g[n] = 1 - A[n]\,\frac{1 - \lfo[n]}{2} \;\in\; [1 - A,\, 1]
\end{equation}
where env $+$ vel is the exponent of eq.~\eqref{eq:cutoff}. Bend glides by a linear ramp over 5\,ms; $w$, $c$, $d_p$, $d_c$ and $A$ over 20\,ms.
A cutoff buffer forces the gliding path (eq.~\eqref{eq:modcutoff} at every 16-sample update); $g$ scales the summed voices.

\subsection{Per-Source Application}

With the note's full-rate ratio $\varrho = 2^{(\mathrm{note}-\mathrm{root})/12} F_{\mathrm{file}}/\Fs$, copy level $\lambda \in \{0,1,2\}$, and $n_0$ a grain chunk's first sample:
\begin{equation}\label{eq:oscstep}
  \Delta\phi[n] = \min\!\left(r[n]\,\frac{f}{\Fs},\; 0.49\right) \ \text{(oscillators, Random)}, \qquad
  \mathrm{inc}[n] = \frac{\varrho\, r[n]}{2^{\lambda}} \ \text{(sample reader)}
\end{equation}
\begin{equation}\label{eq:levels}
  \lambda \leftarrow \lambda + 1 \ \text{ while } \varrho\, r/2^{\lambda} > 2^{3/12}, \qquad
  \lambda \leftarrow \lambda - 1 \ \text{ while } \varrho\, r/2^{\lambda-1} \le 2^{2/12}
\end{equation}
\begin{equation}\label{eq:grainstep}
  s_j' = \bigl\lfloor s_j \cdot r[n_0]/r_j \bigr\rfloor \quad (32.32;\ r_j\text{: the ratio at the grain's start; chunks of } m \le 32 \text{ samples})
\end{equation}
Copy changes crossfade linearly over 2\,ms; a grain step shrinks rather than read past its copy.

\begin{callout}{Bit-Identity Condition}
A buffer reaches the voices as a null pointer when it is neutral for the whole block: pitch $\Leftrightarrow r[n] = 1\ \forall n$;
cutoff $\Leftrightarrow d_c[n]\,\lfo[n] = 0\ \forall n$; the tremolo is skipped while $A = 0$ and not gliding. With null buffers each voice
runs v1.8.0's code path operation for operation, so default sessions render bit-identically to v1.8.0.
\end{callout}

\end{document}
